% Chapter - Modelling, dependency and correlation ............................................................
\chapter{Modelling, dependency, and correlation}
\label{chapter5}

\epigraph{\textit{Numbers have an important story to tell,\\if given a voice.}}{— Florence Nightingale}

The modern scientific conception of nature rests on the idea that observable phenomena can be represented through variables whose relationships follow general mathematical laws. This perspective has deep historical roots and became explicit through developments in analytic geometry and calculus, particularly through the introduction of coordinate systems that allowed geometric relations to be expressed algebraically. The works of Descartes and Newton are commonly regarded as foundational in establishing this mathematical representation of physical reality \cite{descartes_1637_geometry,newton_1687_principia}. 

\medskip

The development of calculus, rooted in the work of Newton and Leibniz, extended this framework by providing tools to describe continuous change. Relationships between quantities could now be expressed as functions \( y = f(x) \), whose rates of change were studied through derivatives and whose accumulated effects were analyzed through integration \cite{leibniz_1684_nova,cauchy_1821_analyse}. 

\medskip

From this period onward, prediction became a defining feature of scientific explanation: if the relation between variables can be expressed mathematically, then under specified assumptions one may compute future or unobserved values. In the presence of uncertainty or stochastic variation, the attempt at prediction links empirical regularities with formal reasoning and lays the conceptual foundations later absorbed into probabilistic modelling \cite{reichenbach_1938_prediction,cox_2006_principles}. 

% The idea of model, explanation and prediction .......................................................................
\section{The idea of model, explanation and prediction}

Scientific investigation proceeds from observation to structured explanation, transforming collections of measurements—\textit{data}—into relations among quantities—\textit{variables}—that can be expressed symbolically and analyzed systematically \cite{hacking_1975_emergence}. 

\medskip 

A classical example is the study of the position of an object, denoted \(x\), as a function of time \(t\). One writes this relation as \(x = f(t)\). In the case of uniform rectilinear motion (constant velocity), the dependence is linear:
\begin{equation}
    x(t) = a + bt ,
    \label{mru_x}
\end{equation}
where \(a\) represents the initial position and \(b\) the constant velocity. The velocity is obtained as the derivative,
\begin{equation}
    \frac{dx}{dt} = b .
    \label{mru_dx}
\end{equation}

\medskip

In the case of uniformly accelerated motion, such as free fall under constant acceleration, the dependence is quadratic:
\begin{equation}
    x(t) = a + bt + ct^{2},
    \label{mrua_x}
\end{equation}
where \(c\) is related to half the acceleration. The velocity, given by differentiation, is linear:
\begin{equation}
    \frac{dx}{dt} = b + 2ct .
    \label{mrua_dx}
\end{equation}

\medskip

A fundamental conceptual distinction in modelling is between \emph{variables} and \emph{parameters}. A variable is a quantity that may take different values across individuals, experiments, or time, whereas a parameter is a \textit{fixed but unknown} constant characterizing the underlying mechanism that generates those variable outcomes. In the examples above, \(t\) is the independent (or explanatory) variable, and \(x(t)\) is the dependent (or response) variable. The constants \(a\), \(b\), and \(c\) are parameters: numerical quantities determining the intercept, slope, and curvature of the function.

\medskip

The distinction between random variables and fixed parameters became explicit in probability theory. For example, in the binomial model, the random variable \(X\) represents the number of observed successes, while \(n\) (the number of trials) and \(p\) (the probability of success) are parameters of the distribution. Historically, Bernoulli treated probabilities as fixed characteristics of repeatable processes \cite{bernoulli_1713_ars}, Laplace incorporated parameters into analytic probability models \cite{laplace_1812_theorie}, and Fisher later clarified in the twentieth century that parameters are fixed constants to be estimated from random samples \cite{fisher_1925_statistical}. 

\medskip

Formally, one denotes by \(Y_1,\dots,Y_n\) random variables representing observed measurements, whose distribution may depend on a parameter \(\theta\), written \(Y_i \sim f(y \mid \theta)\). In the classical frequentist framework, \(\theta\) is unknown but not random; randomness arises from the sampling process, not from the parameter itself.

\medskip

Thus, in a regression model such as 
\[
Y = a + bX + \varepsilon,
\]
the quantities \(X\) and \(Y\) are variables that vary across observations, while \(a\) and \(b\) are parameters describing the systematic structure relating them. The error term \(\varepsilon\) represents stochastic variation around the deterministic component \(a + bX\).

\medskip

For laboratory scientists, this distinction is crucial: measurements such as enzyme activity, concentration, or response time are variables that fluctuate from sample to sample, whereas quantities such as a true mean effect or regression slope are parameters summarizing the stable underlying mechanism that generated those measurements \cite{degroot_2012_probability}. 

% Modelling errors and uncertainty ....................................................................................
\section{Modelling errors and uncertainty}

All empirical measurements exhibit some sort of variability. The first proper formalization of such intuition is normally attributed to Gauss, from who we inherit the Gaussian distribution. In his \textit{"Theoria motus corporum coelestium"} (1809), in the context of astronomical computation and particularly in determining planetary orbits from noisy observational data, he addressed the problem of estimating orbital parameters from imperfect measurements of celestial positions, introducing the method of least squares as a principled way to reconcile multiple inconsistent observations \cite{gauss_1809_theoria}. 

\medskip

His work demonstrated that observational discrepancies could be treated systematically rather than ad hoc, laying the foundation for statistical estimation in both astronomy and later scientific disciplines. Rather than treating discrepancies between theory and observation as mere imperfections, Gauss modeled them explicitly, showing that deviations themselves could be described probabilistically.

\medskip

In modern notation, we express this idea by writing a model in the form
\[
    Y = f(X) + \varepsilon,
\]
where \(f(X)\) represents the systematic and predictable component, and \(\varepsilon\) represents the stochastic deviation from that systematic structure. The symbol \(\varepsilon\) does not denote error in the sense of “mistake,” but rather unobserved influences, natural variability, and measurement imprecision.

\medskip

To make this explicit, one typically assumes that the error term satisfies
\[
    \mathbb{E}[\varepsilon] = 0,
\]
meaning that deviations fluctuate around zero and do not systematically bias the model upward or downward. The variability of the errors is quantified by their variance,
\[
    \mathrm{Var}(\varepsilon) = \sigma^2,
\]
which measures the typical magnitude of fluctuations around the systematic component. This is like saying that the \(\varepsilon\) uncertainty follows a Gaussian distribution.

\medskip

This connects directly with the distinction between population variance and sample variance introduced in earlier chapters. The quantity \(\sigma^2\) represents the population variance of the error distribution—an unknown parameter characterizing intrinsic variability—while its empirical estimate
\[
    s^2 = \frac{1}{n-1}\sum_{i=1}^{n}(x_i-\hat{x}_i)^2
\]
is the sample variance of the residuals, computed from observed data. Thus, uncertainty in modelling is itself summarized through parameters that must be estimated.

\medskip

Kolmogorov’s axiomatization of probability \cite{kolmogorov_1933_grundbegriffe} provided the rigorous mathematical framework for such modelling, embedding \(\varepsilon\) within a probability space and allowing one to treat \(Y\) itself as a random variable whose distribution depends on both systematic structure and stochastic dispersion.

\medskip

Twentieth-century developments reframed error not as failure but as intrinsic uncertainty requiring quantification rather than elimination \cite{cox_2006_principles,mayo_1996_error}. Statistical inference therefore aims not to remove variability, but to measure its magnitude and incorporate it explicitly into conclusions.

\medskip

An important consequence of modelling uncertainty is the principle of propagation of errors. If a quantity \(Z\) depends on measured variables \(X\) and \(Y\), for instance given some function \(Z = g(X,Y)\), then variability in \(X\) and \(Y\) induces variability in \(Z\). Under smoothness conditions, a first-order approximation gives
\[
    \mathrm{Var}(Z) \approx 
    \left(\frac{\partial g}{\partial X}\right)^2 \mathrm{Var}(X)
    +
    \left(\frac{\partial g}{\partial Y}\right)^2 \mathrm{Var}(Y)
    +
    2\frac{\partial g}{\partial X}\frac{\partial g}{\partial Y}\mathrm{Cov}(X,Y),
\]
showing that uncertainty propagates according to the sensitivities (derivatives) of the function.

\medskip

Thus, modelling errors connects calculus, probability, and statistical estimation: derivatives describe how small perturbations influence outcomes, variance quantifies the magnitude of those perturbations, and probabilistic structure allows uncertainty to be carried through mathematical relationships.

\medskip

For experimental researchers, the error term represents biological variability, measurement imprecision, and unobserved influences. The task of properly accounting for it guards against overstated conclusions and irreproducible findings, reinforcing the importance of explicit uncertainty quantification in modern scientific practice \cite{ioannidis_2005_false}. 

% Estimation, inference, evidence .....................................................................................
\section{Estimation, inference, evidence}

Once a model specifies variables and parameters, the central inferential task is to estimate the unknown parameters from observed data. The data are random variables, while the parameters represent fixed but unknown features of the underlying mechanism. Statistical inference seeks principled rules for using the observed sample to learn about those unknown constants.

\medskip

A key conceptual shift introduced by Fisher in the 1920s was the notion of \emph{likelihood} \cite{fisher_1925_statistical}. Suppose that \(Y_1,\dots,Y_n\) are independent observations with density \(f(y \mid \theta)\), where \(\theta\) is an unknown parameter. The joint probability of observing the sample, viewed as a function of \(\theta\), is
\[
L(\theta) = \prod_{i=1}^n f(Y_i \mid \theta),
\]
and is called the likelihood function. Unlike probability, which treats \(\theta\) as fixed and the data as random, likelihood treats the observed data as fixed and considers how plausible different values of \(\theta\) are in light of those data.

\medskip

To illustrate, consider a simple Bernoulli experiment with success probability \(p\). If \(X\) denotes the number of successes in \(n\) trials, then
\[
\mathbb{P}(X=x \mid p) = \binom{n}{x} p^x (1-p)^{n-x}.
\]
Viewed as a function of \(p\), this expression becomes the likelihood \(L(p)\). Maximizing it with respect to \(p\) yields the estimator
\[
\hat{p} = \frac{x}{n},
\]
showing that the maximum likelihood estimator coincides with the intuitive sample proportion.

\medskip

As another example, suppose \(Y_1,\dots,Y_n\) are independent normal random variables with mean \(\mu\) and known variance \(\sigma^2\). The likelihood function is
\[
L(\mu) = \prod_{i=1}^n 
\frac{1}{\sqrt{2\pi\sigma^2}}
\exp\!\left(
-\frac{(Y_i-\mu)^2}{2\sigma^2}
\right).
\]
Maximizing the logarithm of this function with respect to \(\mu\) leads to
\[
\hat{\mu} = \frac{1}{n}\sum_{i=1}^n Y_i,
\]
demonstrating that the sample mean arises naturally as the maximum likelihood estimator.

\medskip

In general, the maximum likelihood estimator is defined as
\[
\hat{\theta} = \arg\max_{\theta} L(\theta),
\]
that is, the parameter value that makes the observed data most plausible under the assumed model. Under regularity conditions, such estimators possess desirable properties: consistency, asymptotic normality, and efficiency, linking estimation to large-sample theory.

\medskip

This framework formalizes the conceptual separation between random variables and fixed parameters: variability arises from the sampling distribution of the estimator \(\hat{\theta}\), not from randomness in \(\theta\) itself within the classical paradigm \cite{degroot_2012_probability}. Thus, uncertainty about parameters is quantified through the variability of estimators across hypothetical repeated samples.

\medskip

Neyman and Pearson extended Fisher’s work by developing formal procedures for testing hypotheses about parameters while controlling long-run error rates \cite{pearson_neyman_1933_efficient}. Their theory introduced Type I and Type II errors and provided systematic criteria for deciding whether observed data constitute sufficient evidence against a specified parameter value.

\medskip

For laboratory scientists, this means that although individual measurements vary unpredictably, the objective of inference is to estimate stable underlying quantities—such as a true mean difference, reaction rate constant, or regression slope—and to quantify the precision of those estimates through confidence intervals and hypothesis tests grounded in sampling theory \cite{cox_2006_principles}. 

% Linear dependence, correlation, and covariance ......................................................................
\section{Linear dependence, correlation, and covariance}

The quantitative study of dependence between variables emerged in the nineteenth century, particularly through the work of Francis Galton on heredity, where he observed that extreme parental traits tended to be followed by less extreme traits in offspring, a phenomenon he called ``regression toward mediocrity.'' His empirical investigations motivated the need for numerical measures of association and systematic methods for fitting linear relationships \cite{galton_1886_regression}.

\medskip

A natural starting point for measuring dependence is \emph{covariance}. Given random variables \(X\) and \(Y\) with means \(\mu_X\) and \(\mu_Y\), the covariance is defined as
\[
    \mathrm{Cov}(X,Y) = \mathbb{E}\big[(X-\mu_X)(Y-\mu_Y)\big].
\]
Covariance measures whether deviations from the mean of one variable tend to be accompanied by deviations of the same or opposite sign in the other. A positive covariance indicates that large values of \(X\) tend to occur with large values of \(Y\), while a negative covariance indicates opposite tendencies.

\medskip

In practice, given data \((x_i,y_i)\), the sample covariance is computed as
\[
    s_{XY} = \frac{1}{n-1}\sum_{i=1}^n (x_i-\bar{x})(y_i-\bar{y}),
\]
providing an empirical estimate of population covariance. However, covariance depends on the measurement units of \(X\) and \(Y\), which makes its magnitude difficult to interpret directly.

\medskip

To address this limitation, Karl Pearson introduced in the 1890s a standardized measure of association now known as the Pearson correlation coefficient \cite{pearson_1895_contributions}. Correlation is defined as covariance divided by the product of standard deviations:
\[
    \rho = \frac{\mathrm{Cov}(X,Y)}{\sigma_X \sigma_Y},
\]
and its sample version is
\[
    r=\frac{\sum (x_i-\bar{x})(y_i-\bar{y})}
    {\sqrt{\sum (x_i-\bar{x})^2}\sqrt{\sum (y_i-\bar{y})^2}}.
\]
Because it is dimensionless, correlation takes values between \(-1\) and \(1\), providing a standardized measure of linear association.

\medskip

While covariance and correlation describe association, regression provides a model for prediction. The method of least squares, introduced systematically by Gauss in 1809 in his work on celestial mechanics \cite{gauss_1809_theoria}, supplied a rule for fitting a straight line to data by minimizing total squared deviations between observed values and the fitted line, transforming empirical approximation into a calculus-based optimization principle.

\medskip

Given observations \((x_i,y_i)\) for \(i=1,\dots,n\), the least squares estimators \(\hat{a},\hat{b}\) minimize 
\[
    S(a,b)=\sum_{i=1}^{n}(y_i-a-bx_i)^2,
\]
and differentiation yields the normal equations
\[
    \frac{\partial S}{\partial a}= -2\sum (y_i-a-bx_i)=0, \qquad \frac{\partial S}{\partial b}= -2\sum x_i(y_i-a-bx_i)=0.
\]

\medskip

Solving these equations gives
\[
    \hat{b}=\frac{\sum (x_i-\bar{x})(y_i-\bar{y})}{\sum (x_i-\bar{x})^2}, \qquad \hat{a}=\bar{y}-\hat{b}\bar{x},
\]
revealing that the regression slope is the ratio of the sample covariance to the sample variance of \(X\). In fact,
\[
    \hat{b}= r\,\frac{s_Y}{s_X},
\]
making explicit the algebraic connection between regression and correlation.

\medskip

A simple linear regression model takes the form
\[
Y = a + bX + \varepsilon,
\]
where \(b\) represents the expected change in \(Y\) associated with a one-unit change in \(X\), and \(\varepsilon\) captures stochastic deviation around the linear trend.

\medskip

This framework generalizes to the linear model \(Y = X\beta + \varepsilon\), unifying regression, analysis of variance, and experimental design within a common algebraic structure \cite{fisher_1925_statistical,fisher_1935_design}. In matrix form, estimation becomes a projection problem in linear algebra.

\medskip

Despite the apparent precision of regression coefficients and correlation values, graphical inspection remains essential, a lesson illustrated by Anscombe in 1973 \cite{anscombe_1973_graphs}, who constructed four datasets with identical means, variances, correlations, and regression lines yet dramatically different scatterplot structures.

\medskip

Anscombe’s quartet demonstrates that numerical summaries alone can conceal curvature, outliers, or clustering, and that statistical modelling must integrate quantitative estimation with visual reasoning to avoid misleading inference.

\medskip

This episode reflects Tukey’s broader advocacy for exploratory data analysis \cite{tukey_1977_eda}, emphasizing that models should be guided by empirical structure rather than imposed mechanically, and reminding researchers that correlation quantifies linear association but does not guarantee linear form, causal interpretation, or robustness to atypical observations.

\medskip

The classical linear model assumes homoscedastic, normally distributed errors—assumptions that simplify estimation and inference but are not universally satisfied in scientific practice—motivating extensions that preserve linear structure while relaxing distributional constraints.

\medskip

A final conceptual distinction is essential: correlation measures statistical association, whereas causation concerns underlying mechanisms. As emphasized in the philosophy of science, statistical regularity alone does not establish that one variable produces changes in another. Salmon argued that scientific explanation requires causal structure rather than mere probabilistic dependence \cite{salmon_1984_scientific}, and Hacking similarly noted that stable statistical patterns may arise without direct causal interaction \cite{hacking_1975_emergence}. Thus, a nonzero correlation—however strong—does not by itself justify causal interpretation; establishing causation requires experimental control, temporal ordering, and theoretical justification beyond statistical association.

% Beyond linearity: GLMs and Interpretation ...........................................................................
\section{Beyond linearity: GLMs and Interpretation}

The classical linear model assumes that the response variable is continuous and normally distributed with constant variance. In many scientific applications, however, outcomes are binary, counts, proportions, or otherwise non-Gaussian. The probabilistic foundation established by Kolmogorov \cite{kolmogorov_1933_grundbegriffe} allows such situations to be treated within a unified framework by specifying appropriate probability distributions rather than forcing normality.

\medskip

Generalized Linear Models (GLMs), systematically developed in the twentieth century, extend linear modelling by separating three components: a random component specifying the distribution of \(Y\), a systematic component given by a linear predictor \(\eta = X\beta\), and a link function \(g\) connecting the mean of the response to the predictor,
\[
g(\mathbb{E}[Y]) = X\beta.
\]
This structure preserves linearity in the parameters while allowing flexibility in the distribution of the response.

\medskip

In most applications, the response distribution belongs to the exponential family, which includes the normal, binomial, and Poisson distributions. This family has mathematical properties that ensure tractable likelihood-based estimation and coherent inferential procedures.

\medskip

For binary outcomes \(Y \in \{0,1\}\), logistic regression assumes
\[
\mathbb{P}(Y=1 \mid X) = \pi(X),
\]
with the logit link
\[
\log\!\left(\frac{\pi(X)}{1-\pi(X)}\right) = X\beta.
\]
Here, coefficients describe changes in log-odds, and exponentiated coefficients represent multiplicative changes in odds.

\medskip

For count data, Poisson regression assumes
\[
Y \sim \text{Poisson}(\lambda),
\qquad
\log(\lambda) = X\beta,
\]
where the logarithmic link ensures positivity of the mean. The coefficients represent multiplicative effects on expected counts, making interpretation natural in epidemiology, ecology, and molecular biology.

\medskip

In practice, real data may exhibit overdispersion, meaning that the variance exceeds the mean in count models. Extensions such as quasi-likelihood methods or negative binomial models address such departures while retaining the linear predictor structure.

\medskip

Estimation in GLMs proceeds via maximum likelihood, linking this framework directly to the inferential principles developed earlier. Under suitable regularity conditions, estimators are consistent and asymptotically normal, permitting construction of confidence intervals and hypothesis tests \cite{cox_2006_principles}.

\medskip

Modern biological analyses, including RNA-seq differential expression, rely heavily on likelihood-based GLM frameworks adapted to count data with dispersion parameters \cite{love_2014_deseq2}. These methods demonstrate how probabilistic modelling, rather than ad hoc transformation, provides principled handling of non-normal outcomes.

\medskip

Finally, although GLMs expand modelling flexibility, statistical association—whether linear or generalized—does not automatically establish causation. Rigorous interpretation requires careful experimental design, control of confounding, and critical scrutiny consistent with broader principles of scientific reasoning \cite{popper_1934_logic,mayo_2018_severe}.
